\explainedfigure{fig_birds_eye_view_fidelity_before_collision}{0.97}{Gap to the other vehicle as drawn in the model's input, in the last steps before two collisions}{For two crashes of a bird's-eye-view model, the true distance between the car and the nearest vehicle (black) is compared with the distance that the 64 by 64 grid shows (red). A hollow circle marks steps at which the vehicle is not in the grid at all. These numbers are copied from the diagnosis note of 10 November (two replayed episodes).}{The oncoming vehicle in the left panel first appears in the grid only 0.9 seconds before impact. In the right panel the grid shows the crossing vehicle 1.7 metres farther away than it really is, and the vehicle disappears from the grid 0.1 seconds before the collision.}{The model's view of the other vehicles is late and coarse: vehicles are drawn as blocks without heading, and the grid ends 17.5 metres ahead. This is a limit of the representation that every model using this grid shares, and it is the leading explanation for the remaining failures.}{fig:fidelity}
